% from ega4-2-fr.tex
\oldpage[IV-2]{82}

\begin{proof}
L'équivalence de a) et b) résulte trivialement de (I, 1.6.5); il est clair
d'autre part que a) entraîne c). Inversement, si c) est vérifiée, on sait que
l'on peut écrire (à identification canonique près)
$N'=\overline N\cap M_{(K')}$. Comme $\overline N$ et $M_{(K')}$ sont par
hypothèse des $A_{(K')}$-modules, il en est de même de $N'$.
\end{proof}

\begin{corollary}[4.8.7]
\label{IV.2.4.8.7-fr}
Sous les hypothèses de (4.8.6), il existe un plus petit sous-corps $K'$ de $K$
contenant $k$ tel que les conditions équivalentes de (4.8.6) soient vérifiées.
\end{corollary}

\begin{proof}
Il suffit de le voir pour la condition c), où cela résulte de Bourbaki,
\emph{Alg.}, chap.~II, 3\textsuperscript{e} éd., \S~8,
n\textsuperscript{o}~6, prop.~6.
\end{proof}

\begin{lemma}[4.8.8]
\label{IV.2.4.8.8-fr}
Les notations étant celles de (4.8.2):
\begin{enumerate}
\item[(i)] Soit
$\overline u:\mathscr F\otimes_kK\to\mathscr G\otimes_kK$ un
$\mathcal O_{X_{(K)}}$-homomorphisme, et soit
$\overline{\mathscr H}\subset
(\mathscr F\otimes_kK)\oplus(\mathscr G\otimes_kK)$ son graphe. Pour que
$\overline u$ soit défini sur $K'$, il faut et il suffit que
$\overline{\mathscr H}$ le soit.
\item[(ii)] Soit $\overline Z$ un sous-préschéma fermé de $X_{(K)}$, et soit
$\overline{\mathscr J}$ l'Idéal quasi-cohérent de $\mathcal O_{X_{(K)}}$
définissant $\overline Z$. Pour que $\overline Z$ soit défini sur $K'$, il
faut et il suffit que $\overline{\mathscr J}$ le soit.
\item[(iii)] Soit $\overline f$ un $K$-morphisme:
$X_{(K)}\to Y_{(K)}$ et soit $\overline Z$ le sous-préschéma de
$X_{(K)}\times_KY_{(K)}$, graphe de $\overline f$ (I, 5.3.11); pour que
$\overline f$ soit défini sur $K'$, il faut et il suffit que $\overline Z$ le
soit.
\end{enumerate}
\end{lemma}

\begin{proof}
L'assertion (ii) résulte aussitôt de (I, 4.4.5). La nécessité dans l'assertion
(iii) est évidente; supposons réciproquement que l'on ait
$\overline Z=Z'_{(K)}$, où $Z'$ est un sous-préschéma de
$(X\times_kY)_{(K')}$; si $p':Z'\to X_{(K')}$ est la restriction à $Z'$ de la
première projection, on sait que $p'_{(K)}$ est un isomorphisme de préschémas,
donc il en est de même de $p'$ (2.7.1, (viii)); mais cela signifie (I, 5.3.11)
que $Z'$ est le graphe d'un $K'$-morphisme
$f':X_{(K')}\to Y_{(K')}$, et l'on a alors
$\overline f=f'_{(K)}$ (I, 5.3.12).

Enfin, pour prouver (i), on peut, en vertu de (4.8.5), supposer que
$X=\operatorname{Spec}(A)$ est affine, $\mathscr F=\widetilde M$,
$\mathscr G=\widetilde N$, où $M$ et $N$ sont des $A$-modules, et
$\overline u=(\overline v)^{\sim}$, où
$\overline v:M_{(K)}\to N_{(K)}$ est un $A_{(K)}$-homomorphisme. Supposons
que le graphe de $\overline v$ soit de la forme $P'_{(K)}$, où $P'$ est un
sous-$A_{(K')}$-module de $M_{(K')}\oplus N_{(K')}$; si
$p':P'\to M_{(K')}$ est la restriction à $P'$ de la première projection, on
sait que $p'\otimes1_K$ est un isomorphisme de $A_{(K)}$-modules, donc, par
fidèle platitude $(0_{\mathrm I},6.4.1)$, $p'$ est un isomorphisme de
$A_{(K')}$-modules, ce qui prouve que $P'$ est le graphe d'un
$A_{(K')}$-homomorphisme $v':M_{(K')}\to N_{(K')}$ et l'on a évidemment
$\overline v=v'\otimes1_{K'}$.
\end{proof}

\begin{proposition}[4.8.9]
\label{IV.2.4.8.9-fr}
Soient $k$ un corps, $K$ une extension de $k$, $X$ un $k$-préschéma,
$\mathscr F$ un $\mathcal O_X$-Module quasi-cohérent. Soit $\overline{\mathscr H}$ un
sous-$\mathcal O_{X_{(K)}}$-Module quasi-cohérent de
$\mathscr F\otimes_kK$. Alors il existe un plus petit corps de définition de
$\overline{\mathscr H}$.
\end{proposition}

\begin{proof}
Lorsque $X$ est affine, l'assertion n'est autre que (4.8.7). Dans le cas
général, considérons un recouvrement $(X_\alpha)$ de $X$ par des ouverts
affines; en vertu de ce qui précède, il existe pour tout $\alpha$ un plus petit
sous-corps $K'_\alpha$ de $K$ contenant $k$ et tel que
$\overline{\mathscr H}|(X_\alpha)_{(K)}$ soit défini sur $K'_\alpha$. En vertu de
(4.8.5), le sous-corps de $K$ engendré par les $K'_\alpha$ est le plus petit
corps de définition de $\overline{\mathscr H}$.
\end{proof}

\begin{corollary}[4.8.10]
\label{IV.2.4.8.10-fr}
Les notations étant celles de (4.8.9), soit $\mathscr G$ un second
$\mathcal O_X$-Module quasi-cohérent, et soit
$\overline u:\mathscr F\otimes_kK\to\mathscr G\otimes_kK$ un
$\mathcal O_{X_{(K)}}$-homomorphisme; alors il existe un plus petit corps de
définition de $\overline u$.
\end{corollary}

